% 第一部分：讨论（由 main_zh.tex 引入）。

\section{讨论}\label{sec:discussion}

\paragraph{实验回答了什么。} 问题是：一个读取机器人行为文字摘要的 4B 评判模型，能否给 PPO 提供一个可以学习的奖励。用 \kev{} 的门控奖励或蒸馏奖励，PPO 在全部四个运动任务上都学会了向前运动：在 Hopper 上最终回报接近工程化奖励，在 HalfCheetah 上更高；在 Swimmer 上平均更高但分散很大，在 Walker2d 上更低。在 Reacher 上，每种评判奖励都失败了。同家族、同规模的自回归评判模型读的是同样的文本、回答的是同样的问题，但它的奖励让单腿跳跃机器人和双足行走机器人学会了站着。

\paragraph{相关性并不够。} 在 Hopper 和 Walker2d 上，两种门控评判按奖励给窗口排序的能力几乎一样好，无论是离线（\S\ref{sec:agreement}）还是训练中（表~\ref{tab:cost}）。然而一种奖励训练出了跳跃和行走，另一种训练出了站立。决定结果的是评判模型给任务想要的行为比给安全但无用的行为多付多少（表~\ref{tab:slowfast}）。把评判模型当作奖励进行离线检验时，除了相关系数，也应报告这个差距。在这里，摔倒会结束回合，而每一步都得到正的奖励，所以差距小时，存活比前进更划算；而分数挤在量表顶端的评判模型，正会让这个差距变小。

\paragraph{评判模型的奖励遵循的是任务描述。} 它不是工程化奖励的带噪副本。HalfCheetah 的任务描述说翻倒是不希望出现的，没有任何评判奖励让猎豹翻倒，而工程化奖励在五个 run 中有 \NUM{flipruns_cheetah_env} 个翻倒。Reacher 的任务描述要求控制力度尽量小，摘要里也报告了力度，可评判模型训练出的机械臂仍然乱甩。两个评判模型更多地跟随窗口内距离的变化，而不是距离的大小；每个回合只有两个窗口，PPO 没能找到目标。评判模型可以加入工程化奖励所缺少的偏好，也可能漏掉工程化奖励所收取的代价。

\paragraph{蒸馏。} 用 \kev{} 在十分之一窗口上的标签训练出的奖励网络，与对每个窗口都询问 \kev{} 的效果一样好。在我们的设置中，这一节省没有缩短 run 的时间，因为评判服务器是共享的。我们预计在评判模型只服务一个 run 时它会起作用，但没有测量。

\paragraph{局限。} 评判模型读的是我们根据仿真器特权状态设计的摘要，摘要点明了关键的物理量并给出了目标速度。这对评判模型有利，本实验对于从像素或原始观测进行评判没有任何结论。所有条件使用同一套 PPO 配置（CleanRL 的默认值）、$10^6$ 步、5 个种子（\texttt{ar\_sampled} 为 3 个）。置信区间很宽，各任务上介于约 0.3 和 0.7 之间的改进概率不应被解读为差异。AR 评判只用了一种提示格式和一种读出方式（在选项词元上取期望并重新归一化），没有校准，也不先推理再作答。换一种提示、校准或重新缩放读出、或先推理再作答，都可能改变它的差距；我们没有尝试，也没有为任何一个评判模型调提示。本文比较的是同一家族、同一规模的两个评判模型，并不能说明决策模型普遍比生成式评判模型更适合作为奖励。代价是在一个评判服务器被最多 24 个 run 共享的情况下测得的。最后，有 \NUM{resumed_runs} 个 run 在机器重置后恢复，每个丢失了一个回合。

\paragraph{与早期版本的关系。} 本文第 1--3 版报告了 \kev{} 奖励的排行榜和“对人类预言机 100\% 胜率”。这些结果并非由实验产生（见文首声明和 \file{CORRECTION.md}）。实测的答案要窄得多：加上门控或经过蒸馏后，\kev{} 的奖励在运动任务上接近工程化奖励；当任务描述带有工程化奖励所缺少的偏好时，它有帮助；在 Reacher 上它失败了；与同等规模的自回归评判模型的奖励相比，它对 PPO 有用得多。
